% ============================================================
% Building Blocks of LLMs -- Part 1
% Section 0: Opening (Slides 1-4)
% Section 1: Linear Algebra (Slides 5-16)
% Section 2: Calculus & Optimization (Slides 17-25)
% ============================================================

% ============================================================
% SECTION 0: OPENING (Slides 1-4)
% ============================================================
\section{Opening}

% ----------------------------------------------------------
% Slide 1: Title Slide
% ----------------------------------------------------------
\begin{frame}[plain]
  \begin{tikzpicture}[remember picture, overlay]
    \fill[aiViolet] (current page.south west) rectangle (current page.north east);
  \end{tikzpicture}

  \vspace{1cm}
  \begin{center}
    {\Huge\bfseries\textcolor{white}{The Building Blocks of LLMs}}\\[0.6em]
    {\Large\textcolor{white!90}{A Mathematical History}}\\[1.5em]
    {\large\textcolor{white!80}{From Archimedes to Attention:}}\\[0.3em]
    {\large\textcolor{white!80}{2,000+ Years of Mathematics Inside Your AI}}\\[2em]
    {\normalsize\textcolor{white!70}{Prof.\ J.\ Osterrieder \hspace{2em} 2026}}
  \end{center}

  \note{
    Welcome slide. Frame the lecture: this is not about hype or ``intelligence.''
    This is about tracing every mathematical component of the LLMs students use daily
    back to the person who invented it.
    Core throughline: ``Every component of an LLM was invented by someone,
    for a specific purpose, often centuries before AI existed.''
    Audience takeaway: ``I already know some of this math ---
    I just didn't know it was inside my AI.''
  }
\end{frame}

% ----------------------------------------------------------
% Slide 2: What's Inside Your AI?
% ----------------------------------------------------------
\begin{frame}{What's Inside Your AI?}

  \begin{columns}[T]
    \begin{column}{0.45\textwidth}
      \vspace{0.5em}
      When you type a question to ChatGPT or Claude\ldots

      \vspace{0.5em}
      The machine that answers is built \alert{entirely from mathematics}.

      \vspace{0.5em}
      \begin{itemize}
        \item Not magic
        \item Not ``intelligence''
        \item \textbf{Specific math}, invented by \textbf{specific people},\\
              at \textbf{specific times}
      \end{itemize}

      \vspace{0.8em}
      \begin{insightbox}
        You already know some of this math.
        You just don't know it's inside your AI yet.
      \end{insightbox}
    \end{column}

    \begin{column}{0.52\textwidth}
      % -- Slide 2: Smartphone with math-topic arrows --
      \begin{tikzpicture}[
        every node/.style={font=\small\bfseries},
        arr/.style={->,>=stealth,thick}
      ]
        % Smartphone outline
        \draw[rounded corners=6pt, thick, fill=gray!10]
          (0,0) rectangle (3,4.5);
        \draw[rounded corners=2pt, fill=white]
          (0.2,0.4) rectangle (2.8,4.1);
        % Chat bubble inside
        \node[draw, rounded corners=3pt, fill=aiViolet!10,
              text width=2cm, align=center, font=\scriptsize]
          (chat) at (1.5,2.5) {AI Response};
        % Six labeled arrows from outside pointing in
        \node[text=linalgTeal] (la) at (-1.8,4.2) {Linear Algebra};
        \draw[arr, linalgTeal] (la) -- (0.4,3.8);
        \node[text=calculusRed] (ca) at (-1.8,3.0) {Calculus};
        \draw[arr, calculusRed] (ca) -- (0.4,3.0);
        \node[text=probBlue] (pr) at (-1.8,1.8) {Probability};
        \draw[arr, probBlue] (pr) -- (0.4,2.2);
        \node[text=infoAmber] (it) at (5.3,4.2) {Info.\ Theory};
        \draw[arr, infoAmber] (it) -- (2.6,3.8);
        \node[text=logicGray] (lo) at (5.3,3.0) {Logic};
        \draw[arr, logicGray] (lo) -- (2.6,3.0);
        \node[text=approxGreen] (fa) at (5.3,1.8) {Func.\ Approx.};
        \draw[arr, approxGreen] (fa) -- (2.6,2.2);
      \end{tikzpicture}
    \end{column}
  \end{columns}

  \note{
    Key message: the AI they use every day is made of math they can learn to understand.
    ``You already know some of this math. You just don't know it's inside your AI yet.
    Today we're going to fix that.''
    Do NOT use hype language. Let the math speak for itself.
  }
\end{frame}

% ----------------------------------------------------------
% Slide 3: The Exploded Diagram -- What We Will Build Today
% ----------------------------------------------------------
\begin{frame}{The Exploded Transformer --- What We Will Build Today}

  \begin{center}
    \visualplaceholder[12cm]{EXPLODED TRANSFORMER DIAGRAM (full-slide, to be generated as figures/exploded-transformer-base.pdf).
      Components to show, color-coded by section:
      (1) Embedding layer -- lookup table, word $\to$ vector [linalgTeal]
      (2) Positional encoding -- sinusoidal waves added to embeddings [seqIndigo]
      (3) Multi-head attention block -- Q/K/V projections, scaled dot-product, concat + linear [linalgTeal + probBlue]
      (4) Add \& Norm (residual connection + layer normalization) [neuronPink]
      (5) Feed-forward network -- two linear layers with activation [approxGreen + neuronPink]
      (6) Second Add \& Norm [neuronPink]
      (7) Repeat $\times N$ layers indicator [logicGray]
      (8) Output linear projection + softmax [probBlue + infoAmber]
      (9) Loss computation -- cross-entropy [infoAmber + calculusRed]
      (10) Backpropagation arrows flowing backward through all layers [calculusRed]
      Layout: vertical stack (input at bottom, output at top), each component a labeled rounded box.
      Intentionally overwhelming on first view.}
  \end{center}

  \vspace{-0.3em}
  {\small\centering\textcolor{gray}{This looks complicated now. By the end, you will know every piece --- and WHO invented it.}\par}

  \note{
    ``This looks complicated right now. Good. By the end of today, you'll know what
    every one of these pieces does and WHO invented it.''
    This slide is meant to overwhelm --- we will return to it in Slide 73
    with full understanding.
    Each color maps to a section of the lecture.
    Do not explain the components yet; just let students absorb the visual.
  }
\end{frame}

% ----------------------------------------------------------
% Slide 4: The Convergence Diagram -- Our Roadmap
% ----------------------------------------------------------
\begin{frame}{Our Roadmap --- The Convergence Diagram}

  \begin{center}
    \visualplaceholder[12cm]{CONVERGENCE DIAGRAM -- INITIAL STATE (to be generated as figures/convergence-0.pdf).
      Layout: central large circle with ``?'' in gray.
      10 satellite nodes in a semicircle (180\textdegree{} arc, evenly spaced), all in gray:
      (1) Linear Algebra [1809--1900s]
      (2) Calculus \& Optimization [c.\,250\,BCE--1986]
      (3) Probability \& Statistics [1654--1933]
      (4) Information Theory [1865--1948]
      (5) Logic \& Computation [1854--1970s]
      (6) Function Approximation [1807--1991]
      (7) The Neuron \& Depth [1943--2016]
      (8) Sequence Modeling [1906--2014]
      (9) The Transformer [2017]
      (10) Scaling to LLMs [2018--2026]
      Gray dashed arrows from each satellite toward center.
      All nodes dim/unlit. Target colors when lit: linalgTeal, calculusRed, probBlue, infoAmber, logicGray, approxGreen, neuronPink, seqIndigo, aiViolet, accentOrange.}
  \end{center}

  \note{
    ``We're going to light up these nodes one by one. Each one is a piece of math
    that somebody invented, often for a completely different purpose, that ended up
    inside the machines you use every day.''
    This diagram will recur at the end of every section, progressively lighting up.
    Emphasize: ten branches of mathematics, developed over 2,000+ years,
    all converge in one architecture.
  }
\end{frame}


% ============================================================
% SECTION 1: LINEAR ALGEBRA -- "The Skeleton of Neural Networks"
% Slides 5-16
% ============================================================
\section{Linear Algebra}

% ----------------------------------------------------------
% Slide 5: Section Divider
% ----------------------------------------------------------
{
\setbeamercolor{background canvas}{bg=digitalBlue}
\begin{frame}[plain]
  \centering
  \vspace{1.5cm}
  \textcolor{white}{\Huge\bfseries Linear Algebra}\\[0.6em]
  \textcolor{white!80}{\large The Skeleton of Neural Networks}\\[1em]
  \textcolor{white!60}{\normalsize 1809--1900s}\\[1.5em]
  \textcolor{white!70}{\itshape``Neural networks are, at their core, matrix multiplication machines.''}

  \note{
    Section transition. Linear algebra provides the fundamental data structures
    (vectors, matrices) and operations (multiplication, dot product) of neural networks.
    Timeline: formal theory begins around 1809 (Grassmann born), but roots in Gauss and earlier.
    Convergence diagram: Section 1 node pulsing.
  }
\end{frame}
}

% ----------------------------------------------------------
% Slide 6: What Is a Vector?
% ----------------------------------------------------------
\begin{frame}{What Is a Vector? --- Forces, Velocities, and Words}

  \begin{columns}[T]
    \begin{column}{0.52\textwidth}
      \begin{personbox}[Hermann Grassmann (1809--1877)]{digitalBlue}
        \small
        \textbf{Origin:} Stettin, Prussia (now Szczecin, Poland)\\
        \textbf{Key work:} \emph{Die lineale Ausdehnungslehre} (1844) ---
        the first general theory of vector spaces in any number of dimensions.\\
        A schoolteacher whose work was so ahead of its time that almost nobody read it.
      \end{personbox}

      \vspace{0.3em}
      \begin{personbox}[Sir William Rowan Hamilton (1805--1865)]{digitalBlue}
        \small
        \textbf{Origin:} Dublin, Ireland\\
        \textbf{Key work:} Invented quaternions (1843) while walking along the Royal Canal
        in Dublin --- famously carved the formula into Brougham Bridge.
      \end{personbox}
    \end{column}

    \begin{column}{0.45\textwidth}
      A \textbf{vector} is a list of numbers that represents something.

      \vspace{0.5em}
      % -- Slide 6: 2D, 3D, and 768-dim vectors --
      \begin{tikzpicture}[scale=0.7]
        % 2D vector
        \draw[->,gray,thin] (-0.3,0)--(2.2,0) node[below,font=\tiny]{$x$};
        \draw[->,gray,thin] (0,-0.3)--(0,1.6) node[left,font=\tiny]{$y$};
        \draw[->,thick,linalgTeal] (0,0)--(1.8,0.9)
          node[right,font=\tiny]{$(2,1)$};
        \node[below,font=\tiny\bfseries] at (1,-.6) {2D};
        % 3D vector (pseudo-isometric)
        \begin{scope}[xshift=3.5cm]
          \draw[->,gray,thin] (-0.3,0)--(1.8,0) node[below,font=\tiny]{$x$};
          \draw[->,gray,thin] (0,-0.3)--(0,1.6) node[left,font=\tiny]{$z$};
          \draw[->,gray,thin] (0,0)--(-0.8,-0.6) node[below,font=\tiny]{$y$};
          \draw[->,thick,linalgTeal] (0,0)--(1.2,1.1);
          \node[below,font=\tiny\bfseries] at (0.6,-.6) {3D};
        \end{scope}
        % 768-dim column vector
        \begin{scope}[xshift=6.8cm, yshift=0.3cm]
          \node[font=\scriptsize,align=center] at (0,1.0) {$\begin{pmatrix}0.23\\-0.87\\0.44\\\vdots\\0.12\end{pmatrix}$};
          \node[below,font=\tiny\bfseries] at (0,-0.7) {768 numbers};
        \end{scope}
      \end{tikzpicture}

      \vspace{0.1em}
      {\scriptsize In an LLM, each word becomes a vector --- a list of 768+ numbers.}

      \vspace{0.3em}
      % -- Slide 6: Brougham Bridge plaque recreation --
      \begin{tikzpicture}
        \node[draw=babylonBrown!70, fill=babylonBrown!8, rounded corners=2pt,
              inner sep=6pt, text width=4.2cm, align=center,
              font=\small] {
          \textsc{Broom Bridge, Dublin}\\[3pt]
          {\itshape Here, on 16th October 1843,}\\
          {\itshape Sir William Rowan Hamilton}\\
          {\itshape carved the formula:}\\[4pt]
          $i^{2} = j^{2} = k^{2} = ijk = -1$
        };
      \end{tikzpicture}
    \end{column}
  \end{columns}

  \note{
    ``Hamilton was so excited that he carved his formula into a bridge. Grassmann's
    far more general idea was ignored for decades. Both were building the language
    that neural networks would eventually speak.''
    LLM connection: every word is converted into a vector (an ``embedding'')
    with 768 or more numbers. Grassmann's abstract vector spaces ARE the embedding
    spaces of modern AI.
    \verifytag{Grassmann's 1844 publication date and characterization of reception.}
  }
\end{frame}

% ----------------------------------------------------------
% Slide 7: Modern Vector Notation -- Gibbs and Heaviside
% ----------------------------------------------------------
\begin{frame}{Modern Vector Notation --- Gibbs and Heaviside}

  \begin{columns}[T]
    \begin{column}{0.50\textwidth}
      \begin{personbox}[Josiah Willard Gibbs (1839--1903)]{digitalBlue}
        \small
        \textbf{Origin:} New Haven, Connecticut, USA\\
        \textbf{Key work:} Extracted the useful vector parts from Hamilton's
        quaternions; created modern vector analysis
        (lecture notes \emph{Elements of Vector Analysis}, c.\,1881--1884).
      \end{personbox}

      \vspace{0.3em}
      \begin{personbox}[Oliver Heaviside (1850--1925)]{digitalBlue}
        \small
        \textbf{Origin:} London, England\\
        \textbf{Key work:} Self-taught electrical engineer who independently developed
        vector calculus; reformulated Maxwell's equations into their compact modern form.
      \end{personbox}
    \end{column}

    \begin{column}{0.47\textwidth}
      % -- Slide 7: Quaternion vs vector notation --
      \begin{tikzpicture}
        % Hamilton box
        \node[draw=digitalBlue!60, fill=digitalBlue!5, rounded corners=2pt,
              text width=2.5cm, align=center, inner sep=5pt,
              font=\scriptsize] (ham) at (0,0) {
          \textbf{Hamilton (1843)}\\[3pt]
          $\mathbf{q}= a + bi + cj + dk$\\[2pt]
          $\mathbf{q}\mathbf{q'} = \ldots$\\
          {\tiny (12 terms)}
        };
        % Gibbs box
        \node[draw=linalgTeal!60, fill=linalgTeal!5, rounded corners=2pt,
              text width=2.5cm, align=center, inner sep=5pt,
              font=\scriptsize] (gib) at (4.2,0) {
          \textbf{Gibbs (1881)}\\[3pt]
          $\vec{a} \cdot \vec{b} = a_1 b_1 + a_2 b_2 + a_3 b_3$\\[2pt]
          {\tiny (3 terms)}
        };
        % Arrow between
        \draw[->,>=stealth,thick,accentOrange] (ham) -- (gib)
          node[midway,above,font=\tiny\bfseries] {simplify};
        % Caption
        \node[font=\scriptsize\itshape, text=gray] at (2.1,-1.3)
          {Same math. Simpler notation.};
      \end{tikzpicture}
    \end{column}
  \end{columns}

  \note{
    ``Gibbs and Heaviside were not trying to build AI. They were trying to make
    physics easier to write down. But their notation is exactly what we use to
    describe neural networks today.''
    \verifytag{Gibbs's lecture notes dated 1881--1884. Confirm exact years of the
    two printings.}
  }
\end{frame}

% ----------------------------------------------------------
% Slide 8: Matrices -- Cayley and Sylvester
% ----------------------------------------------------------
\begin{frame}{Matrices --- Cayley and Sylvester}

  \begin{columns}[T]
    \begin{column}{0.50\textwidth}
      \begin{personbox}[Arthur Cayley (1821--1895)]{digitalBlue}
        \small
        \textbf{Origin:} Richmond, Surrey, England\\
        \textbf{Key work:} ``A Memoir on the Theory of Matrices'' (1858) ---
        established matrices as algebraic objects in their own right.\\
        Also a practicing lawyer for 14 years before returning to academia.
      \end{personbox}

      \vspace{0.3em}
      \begin{personbox}[James Joseph Sylvester (1814--1897)]{digitalBlue}
        \small
        \textbf{Origin:} London, England\\
        \textbf{Key work:} Coined the word ``matrix'' (1850), from Latin for
        ``womb'' --- he saw it as something from which determinants are born.
        Also coined ``discriminant,'' ``invariant,'' and many other terms.
      \end{personbox}
    \end{column}

    \begin{column}{0.47\textwidth}
      % -- Slide 8: Small matrix vs huge LLM matrix --
      \begin{tikzpicture}
        % Small 3x3 matrix
        \node[font=\small, align=center] (small) at (0,1.5) {
          $\begin{pmatrix}
            2 & 0 & 1\\
            3 & 1 & 4\\
            7 & 5 & 2
          \end{pmatrix}$
        };
        \node[font=\tiny\itshape, text=gray, below=1pt of small] {Cayley, 1858};
        % Large LLM weight matrix suggestion
        \node[draw=aiViolet!40, fill=aiViolet!5, minimum width=3.8cm,
              minimum height=1.6cm, rounded corners=1pt] (big) at (0,-1.5) {};
        \node[font=\tiny, text=gray] at (-1.4,-0.9) {0.02};
        \node[font=\tiny, text=gray] at (-0.6,-0.9) {$-$0.7};
        \node[font=\tiny, text=gray] at (0.2,-0.9) {0.31};
        \node[font=\tiny, text=gray] at (0.9,-0.9) {$\cdots$};
        \node[font=\tiny, text=gray] at (-1.4,-1.3) {1.04};
        \node[font=\tiny, text=gray] at (-0.6,-1.3) {0.55};
        \node[font=\tiny, text=gray] at (0.2,-1.3) {$\cdots$};
        \node[font=\tiny, text=gray] at (-0.2,-1.7) {$\vdots$};
        \node[font=\tiny, text=gray] at (0.9,-1.7) {$\ddots$};
        \node[font=\tiny\itshape, text=aiViolet, below=1pt of big]
          {175 billion entries. Same math.};
        % Arrow
        \draw[->,>=stealth,thick,gray!50] (small) -- (big)
          node[midway,right,font=\tiny,text=gray] {scale up};
      \end{tikzpicture}
    \end{column}
  \end{columns}

  \note{
    ``Sylvester literally named it. `Matrix' --- from Latin for `womb.'
    He thought of it as something that gives birth to numbers.
    He had no idea it would give birth to AI.''
    LLM connection: every layer in a neural network is defined by a matrix of
    ``weights.'' When your prompt passes through GPT, it is multiplied by dozens
    of these matrices. Each one was learned during training.
    \verifytag{Cayley's 1858 paper title. Sylvester's coinage of ``matrix''
    in 1850 --- some sources say 1848. Verify exact year and publication.}
  }
\end{frame}

% ----------------------------------------------------------
% Slide 9: Matrix Multiplication -- The Core Operation
% ----------------------------------------------------------
\begin{frame}{Matrix Multiplication --- The Core Operation}

  \begin{columns}[T]
    \begin{column}{0.47\textwidth}
      \begin{personbox}[Carl Friedrich Gauss (1777--1855)]{digitalBlue}
        \small
        \textbf{Origin:} Brunswick, Germany\\
        \textbf{Key work:} Gaussian elimination (c.\,1809--1810) for solving
        linear systems --- laid groundwork for computational linear algebra.
        Originally used for astronomical orbit calculations.
      \end{personbox}

      \vspace{0.5em}
      \begin{insightbox}
        Every neural network layer is a matrix multiplication
        followed by a nonlinear function. That's it.
      \end{insightbox}
    \end{column}

    \begin{column}{0.50\textwidth}
      % -- Slide 9: Matrix multiply + neural net layer --
      \begin{tikzpicture}[scale=0.85]
        % --- TOP: Matrix times vector ---
        \node[font=\scriptsize] at (0,3.8) {\textbf{Matrix $\times$ vector:}};
        \node[font=\scriptsize] (mat) at (0,2.8) {
          $\begin{pmatrix}
            \textcolor{linalgTeal}{w_{11}} & \textcolor{linalgTeal}{w_{12}} & \textcolor{linalgTeal}{w_{13}}\\
            w_{21} & w_{22} & w_{23}\\
            w_{31} & w_{32} & w_{33}
          \end{pmatrix}$};
        \node[font=\scriptsize] at (2.4,2.8) {$\times$};
        \node[font=\scriptsize] at (3.2,2.8) {
          $\begin{pmatrix}x_1\\x_2\\x_3\end{pmatrix}$};
        \node[font=\scriptsize] at (4.2,2.8) {$=$};
        \node[font=\scriptsize] at (5.0,2.8) {
          $\begin{pmatrix}\textcolor{linalgTeal}{y_1}\\y_2\\y_3\end{pmatrix}$};
        % --- BOTTOM: Same as neural net ---
        \node[font=\scriptsize] at (0,1.0) {\textbf{Same operation as a layer:}};
        % Input nodes
        \node[circle, draw, fill=gray!15, inner sep=1.5pt, font=\tiny]
          (i1) at (0,0) {$x_1$};
        \node[circle, draw, fill=gray!15, inner sep=1.5pt, font=\tiny]
          (i2) at (0,-0.7) {$x_2$};
        \node[circle, draw, fill=gray!15, inner sep=1.5pt, font=\tiny]
          (i3) at (0,-1.4) {$x_3$};
        % Output nodes
        \node[circle, draw, fill=linalgTeal!20, inner sep=1.5pt, font=\tiny]
          (o1) at (4,0) {$y_1$};
        \node[circle, draw, fill=linalgTeal!20, inner sep=1.5pt, font=\tiny]
          (o2) at (4,-0.7) {$y_2$};
        \node[circle, draw, fill=linalgTeal!20, inner sep=1.5pt, font=\tiny]
          (o3) at (4,-1.4) {$y_3$};
        % Connections (weights)
        \draw[thin,gray!60] (i1)--(o1); \draw[thin,gray!60] (i1)--(o2); \draw[thin,gray!60] (i1)--(o3);
        \draw[thin,gray!60] (i2)--(o1); \draw[thin,gray!60] (i2)--(o2); \draw[thin,gray!60] (i2)--(o3);
        \draw[thin,gray!60] (i3)--(o1); \draw[thin,gray!60] (i3)--(o2); \draw[thin,gray!60] (i3)--(o3);
        % Highlight one row
        \draw[thick,linalgTeal] (i1)--(o1);
        \draw[thick,linalgTeal] (i2)--(o1);
        \draw[thick,linalgTeal] (i3)--(o1);
        \node[font=\tiny\itshape,text=linalgTeal] at (2,-0.05) {$w_{11},w_{12},w_{13}$};
        % Caption
        \node[font=\scriptsize\itshape,text=gray] at (2.5,-2.1) {Same math. Different picture.};
      \end{tikzpicture}
    \end{column}
  \end{columns}

  \note{
    ``Every time you send a message to an AI, your words pass through about 96 layers.
    Each layer multiplies your data by a matrix. That's roughly 96 matrix multiplications
    just to generate one word of the response.''
    Emphasize: this is NOT an analogy. Neural network layers ARE matrix multiplications.
    The connection is literal and exact.
    \verifytag{Gauss's elimination method typically dated to c.\,1809--1810.
    Method has earlier precedents in Chinese mathematics (Jiuzhang suanshu).}
  }
\end{frame}

% ----------------------------------------------------------
% Slide 10: The Dot Product -- Measuring Similarity
% ----------------------------------------------------------
\begin{frame}{The Dot Product --- Measuring Similarity}

  \begin{columns}[T]
    \begin{column}{0.47\textwidth}
      \textbf{Recipe} (the dot product in plain language):
      \begin{enumerate}
        \item Take two lists of numbers, same length
        \item Multiply matching entries:\\
              first $\times$ first, second $\times$ second, \ldots
        \item Add up all the products
        \item Result is \textbf{one number}:
      \end{enumerate}

      \vspace{0.3em}
      \begin{center}
        \small
        \begin{tabular}{rl}
          Big positive & $=$ very similar \\
          Near zero & $=$ unrelated \\
          Big negative & $=$ opposites
        \end{tabular}
      \end{center}

      \vspace{0.3em}
      {\footnotesize\colorbox{digitalBlue!8}{\parbox{0.9\linewidth}{%
        \textbf{Formal:}\quad $\mathbf{a} \cdot \mathbf{b}
        = a_1 b_1 + a_2 b_2 + \cdots + a_n b_n$}}}
    \end{column}

    \begin{column}{0.50\textwidth}
      % -- Slide 10: Dot product -- three configurations --
      \begin{tikzpicture}[scale=0.75]
        % (a) Same direction
        \begin{scope}[xshift=0cm]
          \draw[->,thick,linalgTeal] (0,0) -- (1.5,0.6);
          \draw[->,thick,accentOrange] (0,0) -- (1.6,0.4);
          \node[font=\tiny\bfseries,text=linalgTeal] at (0.8,-0.5) {large $+$};
          \node[font=\tiny] at (0.8,-0.9) {(a) similar};
        \end{scope}
        % (b) Perpendicular
        \begin{scope}[xshift=3cm]
          \draw[->,thick,linalgTeal] (0,0) -- (1.5,0);
          \draw[->,thick,accentOrange] (0,0) -- (0,1.2);
          \node[font=\tiny\bfseries,text=gray] at (0.8,-0.5) {$\approx 0$};
          \node[font=\tiny] at (0.8,-0.9) {(b) unrelated};
        \end{scope}
        % (c) Opposite
        \begin{scope}[xshift=6cm]
          \draw[->,thick,linalgTeal] (0,0) -- (1.5,0.3);
          \draw[->,thick,accentOrange] (0,0) -- (-1.3,-0.3);
          \node[font=\tiny\bfseries,text=calculusRed] at (0.2,-0.5) {large $-$};
          \node[font=\tiny] at (0.2,-0.9) {(c) opposite};
        \end{scope}
      \end{tikzpicture}

      \vspace{0.2em}
      {\scriptsize\itshape The dot product measures how much two vectors agree.}

      \vspace{0.3em}
      % -- Slide 10: Attention dot-product "it" -> "cat" --
      \begin{tikzpicture}
        \node[draw, rounded corners=2pt, fill=probBlue!10, inner sep=4pt,
              font=\small\bfseries] (it) at (0,0) {``it''};
        \node[draw, rounded corners=2pt, fill=probBlue!10, inner sep=4pt,
              font=\small\bfseries] (cat) at (4,0) {``cat''};
        \draw[->,>=stealth, line width=2.5pt, probBlue] (it) -- (cat)
          node[midway, above, font=\tiny] {large dot product};
        \node[font=\scriptsize\itshape, text=gray, text width=4.5cm, align=center]
          at (2,-0.9) {The model ``attends'' \textbf{it} to \textbf{cat}.\\
          This is how it resolves what ``it'' refers to.};
      \end{tikzpicture}
    \end{column}
  \end{columns}

  \note{
    ``The dot product is the simplest possible measure of similarity between two vectors.
    Multiply matching entries and add up. That's it. This tiny operation, repeated
    billions of times, is what lets an LLM understand which words relate to which.''
    Historical: formalized by Gibbs (1880s), implicit in earlier work by Grassmann.
    LLM connection: in the attention mechanism, the dot product of a Query vector and
    a Key vector determines how much ``attention'' one word pays to another.
    Every attention score in every LLM is a dot product.
  }
\end{frame}

% ----------------------------------------------------------
% Slide 11: Cosine Similarity
% ----------------------------------------------------------
\begin{frame}{Cosine Similarity --- Dot Product, Normalized}

  \begin{columns}[T]
    \begin{column}{0.45\textwidth}
      \textbf{Recipe:}
      \begin{enumerate}
        \item Compute the dot product
        \item Divide by the length of each vector
        \item Result is always between $-1$ and $+1$
      \end{enumerate}

      \vspace{0.3em}
      Think of it as: \emph{how much do these two vectors point the same way?}
      (Ignoring how long they are.)

      \vspace{0.5em}
      {\footnotesize\colorbox{digitalBlue!8}{\parbox{0.9\linewidth}{%
        \textbf{Formal:}\quad $\cos(\theta) = \dfrac{\mathbf{a} \cdot \mathbf{b}}%
        {\lvert\mathbf{a}\rvert \;\lvert\mathbf{b}\rvert}$}}}

      \vspace{0.5em}
      \begin{insightbox}
        When you ask an LLM to find ``similar words,'' it computes cosine similarity.
        The math is from the 1880s. The application is from last Tuesday.
      \end{insightbox}
    \end{column}

    \begin{column}{0.52\textwidth}
      % -- Slide 11: Cosine similarity in embedding space --
      \begin{tikzpicture}[scale=0.9]
        % Axes
        \draw[->,gray,thin] (-0.3,0)--(4.5,0);
        \draw[->,gray,thin] (0,-0.3)--(0,3.5);
        % King vector
        \draw[->,thick,linalgTeal] (0,0) -- (3,2.6)
          node[right,font=\scriptsize\bfseries] {king};
        % Queen vector (close to king)
        \draw[->,thick,probBlue] (0,0) -- (2.5,2.9)
          node[above,font=\scriptsize\bfseries] {queen};
        % Banana vector (far away)
        \draw[->,thick,accentOrange] (0,0) -- (-0.3,2.2)
          node[left,font=\scriptsize\bfseries] {banana};
        % Small angle arc between king and queen
        \draw[linalgTeal, thick] (1.0,0.87) arc[start angle=41, end angle=49, radius=1.33];
        \node[font=\tiny,text=linalgTeal] at (1.8,1.9) {small $\theta$};
        % Large angle arc between king and banana
        \draw[calculusRed, thick] (0.3,1.0) arc[start angle=73, end angle=41, radius=1.05];
        \node[font=\tiny,text=calculusRed] at (1.0,1.5) {large $\theta$};
        % Caption
        \node[font=\scriptsize\itshape, text=gray, text width=4cm, align=center]
          at (2.2,-0.7) {Cosine similarity = the angle between word vectors.};
      \end{tikzpicture}
    \end{column}
  \end{columns}

  \note{
    ``When you ask an LLM to find `similar words,' it computes cosine similarity
    in embedding space. The math is from the 1880s. The application is from
    last Tuesday.''
    This is how search engines, recommendation systems, and LLMs measure
    semantic similarity between pieces of text.
  }
\end{frame}

% ----------------------------------------------------------
% Slide 12: Eigenvalues and Eigenvectors [EXTENDED]
% ----------------------------------------------------------
\begin{frame}{Eigenvalues --- The Natural Modes \hfill {\small\textcolor{gray}{[EXTENDED]}}}

  \begin{columns}[T]
    \begin{column}{0.50\textwidth}
      \begin{personbox}[Leonhard Euler (1707--1783)]{digitalBlue}
        \small
        \textbf{Origin:} Basel, Switzerland / St.\ Petersburg / Berlin\\
        \textbf{Key work:} First encountered eigenvalue-like concepts studying
        vibrating strings and rotating bodies (1740s--1750s).
      \end{personbox}

      \vspace{0.2em}
      \begin{personbox}[Augustin-Louis Cauchy (1789--1857)]{digitalBlue}
        \small
        \textbf{Origin:} Paris, France\\
        \textbf{Key work:} Proved symmetric matrices have real eigenvalues (1829).
      \end{personbox}

      \vspace{0.2em}
      \begin{personbox}[David Hilbert (1862--1943)]{digitalBlue}
        \small
        \textbf{Origin:} K\"onigsberg / G\"ottingen, Germany\\
        \textbf{Key work:} Extended eigenvalue theory to infinite-dimensional
        spaces (spectral theory, early 1900s).
      \end{personbox}
    \end{column}

    \begin{column}{0.47\textwidth}
      % -- Slide 12: Eigenvector visualization --
      \begin{tikzpicture}[scale=0.75]
        % Before transformation (left)
        \node[font=\scriptsize\bfseries] at (1.2,2.8) {Before};
        \draw[->,gray,thin] (-0.3,0)--(2.5,0);
        \draw[->,gray,thin] (0,-0.3)--(0,2.5);
        \draw[->,thin,gray!50] (0,0)--(1.5,1.5);
        \draw[->,thin,gray!50] (0,0)--(0.5,1.8);
        \draw[->,thin,gray!50] (0,0)--(1.8,0.5);
        \draw[->,thick,linalgTeal] (0,0)--(2.0,0)
          node[below,font=\tiny] {$\vec{e}_1$};
        \draw[->,thick,accentOrange] (0,0)--(0,1.8)
          node[left,font=\tiny] {$\vec{e}_2$};
        % After transformation (right)
        \begin{scope}[xshift=4.5cm]
          \node[font=\scriptsize\bfseries] at (1.5,2.8) {After $A\vec{v}$};
          \draw[->,gray,thin] (-0.5,0)--(3.0,0);
          \draw[->,gray,thin] (0,-0.5)--(0,2.5);
          % Rotated arrows (non-eigenvectors)
          \draw[->,thin,gray!50] (0,0)--(2.0,0.8);
          \draw[->,thin,gray!50] (0,0)--(1.2,1.5);
          \draw[->,thin,gray!50] (0,0)--(2.2,-0.2);
          % Eigenvectors: only stretched, not rotated
          \draw[->,very thick,linalgTeal] (0,0)--(2.6,0)
            node[below,font=\tiny] {$\lambda_1 \vec{e}_1$};
          \draw[->,very thick,accentOrange] (0,0)--(0,1.0)
            node[left,font=\tiny] {$\lambda_2 \vec{e}_2$};
        \end{scope}
        % Caption
        \node[font=\scriptsize\itshape, text=gray, text width=6cm, align=center]
          at (3.5,-0.9) {Eigenvectors: only stretch or shrink, never rotate.};
      \end{tikzpicture}

      \vspace{0.3em}
      {\small\textbf{LLM connection:} Principal Component Analysis (PCA)
        uses eigenvalues to visualize and understand what LLMs learn.
        ``What does this neuron represent?'' --- eigenvalue analysis helps answer.}
    \end{column}
  \end{columns}

  \note{
    This is an [EXTENDED] slide --- can be cut if short on time.
    Eigenvalues reveal the ``natural modes'' of a transformation.
    They appear in dimensionality reduction (PCA), used for understanding
    and compressing neural network representations.
    \verifytag{Euler's eigenvalue work --- verify precise context
    (vibrating membranes vs.\ rotating bodies) and dates.
    Cauchy's 1829 result on symmetric matrices.}
  }
\end{frame}

% ----------------------------------------------------------
% Slide 13: High-Dimensional Spaces -- Where Words Live
% ----------------------------------------------------------
\begin{frame}{High-Dimensional Spaces --- Where Words Live}

  \begin{columns}[T]
    \begin{column}{0.48\textwidth}
      LLMs represent words as vectors in spaces with
      \textbf{hundreds or thousands of dimensions}.

      \vspace{0.3em}
      The math works the same as in 2D or 3D --- we just can't draw it.

      \vspace{0.5em}
      \begin{personbox}[Richard Bellman (1920--1984)]{digitalBlue}
        \small
        \textbf{Origin:} New York City, USA\\
        \textbf{Key work:} Coined ``the curse of dimensionality'' (c.\,1957--1961):
        in high dimensions, data becomes sparse and distances behave
        counterintuitively.
      \end{personbox}
    \end{column}

    \begin{column}{0.49\textwidth}
      % -- Slide 13: Word vector parallelogram --
      \begin{tikzpicture}[scale=0.85]
        % Word points
        \node[circle, fill=linalgTeal, inner sep=2pt, label={[font=\scriptsize]right:king}]
          (king) at (3.5,2.5) {};
        \node[circle, fill=linalgTeal, inner sep=2pt, label={[font=\scriptsize]right:queen}]
          (queen) at (3.8,3.8) {};
        \node[circle, fill=probBlue, inner sep=2pt, label={[font=\scriptsize]left:man}]
          (man) at (1.0,1.2) {};
        \node[circle, fill=probBlue, inner sep=2pt, label={[font=\scriptsize]left:woman}]
          (woman) at (1.3,2.5) {};
        % Extra words
        \node[circle, fill=gray!40, inner sep=1.5pt, label={[font=\tiny]right:dog}]
          at (4.5,0.8) {};
        \node[circle, fill=gray!40, inner sep=1.5pt, label={[font=\tiny]right:cat}]
          at (4.2,1.2) {};
        % Parallelogram arrows
        \draw[->,thick,accentOrange] (man) -- (king)
          node[midway,below right,font=\tiny] {royalty};
        \draw[->,thick,accentOrange] (woman) -- (queen);
        \draw[->,dashed,gray] (man) -- (woman)
          node[midway,left,font=\tiny] {gender};
        \draw[->,dashed,gray] (king) -- (queen);
        % Caption
        \node[font=\scriptsize\itshape, text=gray, text width=5cm, align=center]
          at (2.5,-0.3) {Projected from 768 dims to 2.\\Mikolov et al., 2013};
      \end{tikzpicture}

      \vspace{0.3em}
      {\small\textbf{Tomas Mikolov} et al.\ demonstrated this with
        \textbf{Word2Vec} (2013) --- word vectors trained on large text corpora
        capture semantic relationships as \emph{linear directions}.}
    \end{column}
  \end{columns}

  \note{
    ``In 768 dimensions, words that mean similar things end up near each other.
    `Dog' is near `puppy.' `Paris' is near `France.' The AI doesn't know what
    these words mean --- it only knows their coordinates. But the coordinates
    capture meaning. Mikolov showed this in 2013 and it stunned the field.''
    \verifytag{Confirm both Mikolov et al.\ 2013 papers: ``Efficient Estimation
    of Word Representations in Vector Space'' (arXiv) and ``Distributed
    Representations of Words and Phrases'' (NeurIPS).}
    \verifytag{Bellman coined ``curse of dimensionality'' --- may be from
    his 1961 book \emph{Adaptive Control Processes}, not 1957. Verify.}
  }
\end{frame}

% ----------------------------------------------------------
% INTERACTIVE MOMENT 1 (after Slide 13)
% ----------------------------------------------------------
\begin{frame}{Quick Experiment --- Word Vector Arithmetic}

  \begin{center}
    {\Large\bfseries\textcolor{aiViolet}{What does this equal?}}

    \vspace{1.5em}
    {\Huge $\underset{\text{king}}{\mathbf{v}} \;-\;
           \underset{\text{man}}{\mathbf{v}} \;+\;
           \underset{\text{woman}}{\mathbf{v}} \;=\; ?$}

    \vspace{1.5em}
    % -- Slide 13i: Reveal -- parallelogram + QUEEN --
    \begin{tikzpicture}[scale=0.9]
      % Parallelogram
      \node[circle, fill=linalgTeal, inner sep=2.5pt,
            label={[font=\small]below left:man}]
        (man) at (0,0) {};
      \node[circle, fill=linalgTeal, inner sep=2.5pt,
            label={[font=\small]above left:woman}]
        (woman) at (0.4,2.0) {};
      \node[circle, fill=accentOrange, inner sep=2.5pt,
            label={[font=\small]below right:king}]
        (king) at (3.5,0.3) {};
      \node[circle, fill=accentOrange, inner sep=3pt,
            label={[font=\large\bfseries,text=accentOrange]above right:QUEEN \checkmark}]
        (queen) at (3.9,2.3) {};
      % Arrows
      \draw[->,thick,gray] (man) -- (king);
      \draw[->,thick,gray] (woman) -- (queen);
      \draw[->,thick,probBlue] (man) -- (woman);
      \draw[->,thick,probBlue] (king) -- (queen);
      % Dashed analogy arrow
      \draw[->,dashed,thick,calculusRed] (king) -- (queen)
        node[midway,right,font=\tiny]{$-$man $+$ woman};
      % Attribution
      \node[font=\scriptsize\itshape, text=gray] at (2,-1.0)
        {Mikolov et al., Word2Vec, 2013};
    \end{tikzpicture}
  \end{center}

  \note{
    INTERACTIVE MOMENT 1 (approx.\ 2 minutes).
    Ask students: ``If we take the vector for `king,' subtract the vector for `man,'
    and add the vector for `woman' --- what word's vector do we land nearest to?''
    Let them guess. Collect answers.
    Reveal: ``queen.'' This works because of the geometry of high-dimensional vector
    spaces, as Mikolov's Word2Vec demonstrated in 2013.
    This is not a trick --- it is a mathematical property of the embedding space.
    Grassmann's abstract vector spaces from 1844 make this possible.
  }
\end{frame}

% ----------------------------------------------------------
% Slide 14: From Words to Numbers -- The Embedding
% ----------------------------------------------------------
\begin{frame}{From Words to Numbers --- The Embedding}

  \begin{center}
    % -- Slide 14: Embedding pipeline --
    \begin{tikzpicture}[
      bx/.style={draw, rounded corners=3pt, inner sep=4pt,
                 font=\scriptsize, align=center, minimum height=0.8cm},
      arr/.style={->,>=stealth,thick,gray}
    ]
      % Step 1: Sentence
      \node[bx, fill=gray!10, text width=2.8cm] (sent) at (0,0)
        {``The cat sat\\on the mat''};
      \node[font=\tiny\bfseries, above=1pt of sent] {Sentence};
      % Step 2: Tokens
      \node[bx, fill=linalgTeal!10, text width=2.2cm] (tok) at (4,0)
        {The\,|\,cat\,|\,sat\\on\,|\,the\,|\,mat};
      \node[font=\tiny\bfseries, above=1pt of tok] {Tokens};
      % Step 3: Vectors
      \node[bx, fill=linalgTeal!20, text width=2.2cm] (vec) at (7.8,0)
        {$\begin{pmatrix}0.23\\-0.87\\\vdots\end{pmatrix} \cdots$\\{\tiny 768 numbers each}};
      \node[font=\tiny\bfseries, above=1pt of vec] {Vectors};
      % Step 4: Matrix
      \node[bx, fill=aiViolet!15, text width=2.0cm] (mat) at (11.2,0)
        {$6 \times 768$\\matrix};
      \node[font=\tiny\bfseries, above=1pt of mat] {Matrix};
      % Arrows
      \draw[arr] (sent) -- (tok) node[midway,above,font=\tiny] {tokenize};
      \draw[arr] (tok) -- (vec) node[midway,above,font=\tiny] {embed};
      \draw[arr] (vec) -- (mat) node[midway,above,font=\tiny] {stack};
    \end{tikzpicture}

    \vspace{0.3em}
    {\small\bfseries\textcolor{aiViolet}{Your sentence is now a matrix.}
    Everything that follows is linear algebra.}
  \end{center}

  \vspace{0.3em}
  {\small\textbf{Historical lineage:}
    Grassmann (1844, abstract vector spaces)
    $\to$ Mikolov (2013, Word2Vec learned embeddings)
    $\to$ Modern LLM embedding layers (2017--present, learned jointly with the transformer)}

  \note{
    ``Grassmann imagined abstract spaces with any number of dimensions. Cayley gave
    us matrices. Mikolov showed words could live in those spaces. Now, 170 years after
    Grassmann, your words live in his spaces.''
    Key point: after embedding, the sentence IS a matrix. Every operation that follows
    is a matrix operation. This is the bridge from language to linear algebra.
  }
\end{frame}

% ----------------------------------------------------------
% Slide 15: Attention as Matrix Multiplication (Preview)
% ----------------------------------------------------------
\begin{frame}{Attention as Matrix Multiplication --- A Preview}

  \begin{columns}[T]
    \begin{column}{0.55\textwidth}
      \textbf{The attention recipe} (conceptual pipeline):

      \vspace{0.3em}
      \begin{enumerate}
        \item Start with input matrix $X$ (sentence as matrix)
        \item Multiply by $W_Q$ $\to$ get Queries\\
              {\scriptsize (matrix multiplication --- Cayley)}
        \item Multiply by $W_K$ $\to$ get Keys\\
              {\scriptsize (matrix multiplication --- Cayley)}
        \item Multiply by $W_V$ $\to$ get Values\\
              {\scriptsize (matrix multiplication --- Cayley)}
        \item Scores $=$ dot product of each Query with each Key\\
              {\scriptsize (dot product --- Gibbs, Slide 10)}
        \item Normalize scores with softmax\\
              {\scriptsize (probability --- Section 3)}
        \item Output $=$ scores $\times$ Values
      \end{enumerate}

      \vspace{0.3em}
      {\footnotesize\colorbox{digitalBlue!8}{\parbox{0.92\linewidth}{%
        \textbf{Formal (preview):}\quad
        $\text{Attn} = \text{softmax}\!\left(\frac{QK^T}{\sqrt{d}}\right) V$\\[2pt]
        $K^T$ $=$ ``flip the matrix on its side'' (transpose).
        $\sqrt{d}$ $=$ scaling factor.\\
        We will unpack this fully in Section~9.}}}
    \end{column}

    \begin{column}{0.42\textwidth}
      % -- Slide 15: Attention Q,K,V pipeline --
      \begin{tikzpicture}[
        bx/.style={draw, rounded corners=2pt, inner sep=3pt,
                   font=\tiny, minimum width=0.7cm, minimum height=0.5cm},
        arr/.style={->,>=stealth,thin}
      ]
        % Input X
        \node[bx, fill=gray!15, font=\tiny\bfseries] (X) at (0,1.5) {$X$};
        % Three parallel paths
        \node[bx, fill=linalgTeal!15] (wq) at (1.5,2.8) {$W_Q$};
        \node[bx, fill=linalgTeal!15] (wk) at (1.5,1.5) {$W_K$};
        \node[bx, fill=linalgTeal!15] (wv) at (1.5,0.2) {$W_V$};
        \node[bx, fill=linalgTeal!25] (Q) at (3.0,2.8) {$Q$};
        \node[bx, fill=linalgTeal!25] (K) at (3.0,1.5) {$K$};
        \node[bx, fill=linalgTeal!25] (V) at (3.0,0.2) {$V$};
        % Arrows: X to W matrices
        \draw[arr] (X) |- (wq);
        \draw[arr] (X) -- (wk);
        \draw[arr] (X) |- (wv);
        % Arrows: W to Q,K,V
        \draw[arr] (wq) -- (Q);
        \draw[arr] (wk) -- (K);
        \draw[arr] (wv) -- (V);
        % Dot product
        \node[bx, fill=probBlue!15] (scores) at (4.5,2.1) {\scriptsize $QK^T$};
        \draw[arr] (Q) -| (scores);
        \draw[arr] (K) -| (scores);
        % Softmax
        \node[bx, fill=accentOrange!15] (sm) at (4.5,0.9) {\tiny softmax};
        \draw[arr] (scores) -- (sm);
        % Times V
        \node[bx, fill=aiViolet!15, font=\tiny\bfseries] (out) at (4.5,-0.2) {Out};
        \draw[arr] (sm) -- (out);
        \draw[arr] (V) -- (out);
      \end{tikzpicture}
    \end{column}
  \end{columns}

  \note{
    ``The `attention' mechanism --- the thing that lets an AI understand context ---
    is three matrix multiplications, a dot product, and a softmax. That's it.
    The math was all invented by the 1880s. The architecture was invented in 2017.''
    This is a PREVIEW. We return to attention in detail in Section 9.
    The point: show students early that ``attention'' is not magic --- it is the
    linear algebra they just learned.
    Do NOT let students get stuck on the formal equation. The conceptual pipeline
    is what matters here.
  }
\end{frame}

% ----------------------------------------------------------
% Slide 16: Convergence Diagram Update -- Linear Algebra
% ----------------------------------------------------------
\begin{frame}{Building the Diagram --- Linear Algebra Complete}

  \begin{columns}[T]
    \begin{column}{0.55\textwidth}
      \begin{center}
        \visualplaceholder[6.5cm]{CONVERGENCE DIAGRAM -- VARIANT 1 (figures/convergence-1.pdf).
          Node (1) ``Linear Algebra'' lit in linalgTeal with solid arrow to center.
          Label: ``Matrix multiplication = neural network layers.''
          Nodes (2)--(10) remain gray with dashed arrows.
          Central node still ``?'' in gray.}
      \end{center}
    \end{column}

    \begin{column}{0.42\textwidth}
      \textbf{Summary --- who built the skeleton:}

      \vspace{0.3em}
      {\small
      \begin{tabular}{@{}rl@{}}
        1843 & Hamilton --- quaternions \\
        1844 & Grassmann --- vector spaces \\
        1850 & Sylvester --- named ``matrix'' \\
        1858 & Cayley --- matrix theory \\
        1880s & Gibbs/Heaviside --- notation \\
        1809+ & Gauss --- elimination \\
        2013 & Mikolov --- Word2Vec \\
      \end{tabular}}

      \vspace{0.5em}
      \begin{insightbox}
        Every layer of an LLM is a matrix multiplication.
        Every attention score is a dot product.
        The math of 1850 is the engine of 2026.
      \end{insightbox}
    \end{column}
  \end{columns}

  \note{
    Transition: ``Now we have the skeleton. But a skeleton doesn't move. For that,
    we need calculus.''
    Recap all six mathematicians and Mikolov. Emphasize the time span:
    from Grassmann's ignored 1844 paper to the core of every modern LLM.
    Connection arrow: ``Every layer of an LLM is a matrix multiplication.
    Every attention score is a dot product. The math of 1850 is the engine of 2026.''
  }
\end{frame}


% ============================================================
% SECTION 2: CALCULUS & OPTIMIZATION -- "How Neural Networks Learn"
% Slides 17-25
% ============================================================
\section{Calculus \& Optimization}

% ----------------------------------------------------------
% Slide 17: Section Divider
% ----------------------------------------------------------
{
\setbeamercolor{background canvas}{bg=enlightenmentTeal}
\begin{frame}[plain]
  \centering
  \vspace{1.5cm}
  \textcolor{white}{\Huge\bfseries Calculus \& Optimization}\\[0.6em]
  \textcolor{white!80}{\large How Neural Networks Learn}\\[1em]
  \textcolor{white!60}{\normalsize c.\,250 BCE -- 1986}\\[1.5em]
  \textcolor{white!70}{\itshape``A neural network learns by rolling downhill
    on a landscape made of calculus.''}

  \note{
    Section transition. Without calculus, you can BUILD a neural network but you
    cannot TRAIN it. Calculus provides the learning algorithm.
    Timeline: from Archimedes' proto-calculus (c.\,250 BCE) to Rumelhart, Hinton,
    and Williams' backpropagation paper (1986).
    Convergence diagram: Section 2 node pulsing.
  }
\end{frame}
}

% ----------------------------------------------------------
% Slide 18: Roots in Antiquity -- Archimedes
% ----------------------------------------------------------
\begin{frame}{Roots in Antiquity --- Archimedes and Infinitesimals}

  \begin{columns}[T]
    \begin{column}{0.50\textwidth}
      \begin{personbox}[Archimedes of Syracuse]{enlightenmentTeal}
        \small
        \textbf{Origin:} Syracuse, Sicily\\
        \textbf{Key work:} The ``method of exhaustion'' --- calculated areas and
        volumes by approximating with ever-finer polygons.\\[3pt]
        His lost \emph{Method} manuscript, rediscovered in 1906 in the
        Archimedes Palimpsest, shows he used infinitesimal reasoning
        more explicitly than previously known.
      \end{personbox}

      \vspace{0.5em}
      \begin{insightbox}
        The core idea of calculus: approximate with simple pieces, take the limit.
        Archimedes had this 1,900 years before Newton.
      \end{insightbox}
    \end{column}

    \begin{column}{0.47\textwidth}
      % -- Slide 18: Method of exhaustion --
      \begin{tikzpicture}[scale=0.65]
        % Circle 1: inscribed square
        \begin{scope}[xshift=0cm]
          \draw[thick,gray] (0,0) circle(1.2);
          \draw[thick,calculusRed] (1.2,0)--(0,1.2)--(-1.2,0)--(0,-1.2)--cycle;
          \node[font=\tiny,below] at (0,-1.5) {4 sides};
        \end{scope}
        % Circle 2: inscribed octagon
        \begin{scope}[xshift=3cm]
          \draw[thick,gray] (0,0) circle(1.2);
          \draw[thick,calculusRed]
            (1.2,0)--(0.85,0.85)--(0,1.2)--(-0.85,0.85)
            --(-1.2,0)--(-0.85,-0.85)--(0,-1.2)--(0.85,-0.85)--cycle;
          \node[font=\tiny,below] at (0,-1.5) {8 sides};
        \end{scope}
        % Circle 3: inscribed 16-gon (approximate with smoother polygon)
        \begin{scope}[xshift=6cm]
          \draw[thick,gray] (0,0) circle(1.2);
          \draw[thick,calculusRed]
            (1.2,0)--(1.11,0.46)--(0.85,0.85)--(0.46,1.11)
            --(0,1.2)--(-0.46,1.11)--(-0.85,0.85)--(-1.11,0.46)
            --(-1.2,0)--(-1.11,-0.46)--(-0.85,-0.85)--(-0.46,-1.11)
            --(0,-1.2)--(0.46,-1.11)--(0.85,-0.85)--(1.11,-0.46)--cycle;
          \node[font=\tiny,below] at (0,-1.5) {16 sides};
        \end{scope}
        % Ellipsis
        \node[font=\large] at (8.3,0) {$\cdots$};
        % Arrow
        \node[font=\scriptsize,right] at (9,0) {$= \pi r^2$};
        % Caption
        \node[font=\scriptsize\itshape, text=gray, text width=5cm, align=center]
          at (4.5,-2.3) {Approximate, refine, repeat.\\This is what calculus does.};
      \end{tikzpicture}
    \end{column}
  \end{columns}

  \note{
    ``Archimedes was doing proto-calculus in the 3rd century BCE,
    1,900 years before Newton.''
    The Archimedes Palimpsest was rediscovered in 1906 --- a medieval prayer book
    written over his mathematical manuscripts. Modern imaging technology revealed
    the hidden text, showing more sophisticated infinitesimal reasoning than
    historians had expected.
  }
\end{frame}

% ----------------------------------------------------------
% Slide 19: The Derivative -- Measuring Change
% ----------------------------------------------------------
\begin{frame}{The Derivative --- Measuring Change}

  \begin{columns}[T]
    \begin{column}{0.48\textwidth}
      \begin{personbox}[Pierre de Fermat (1601--1665)]{enlightenmentTeal}
        \small
        \textbf{Origin:} Beaumont-de-Lomagne, France\\
        \textbf{Key work:} Method for tangent lines (c.\,1636--1638);
        found maxima and minima by setting a difference quotient to zero.
      \end{personbox}

      \vspace{0.2em}
      \begin{personbox}[Isaac Newton (1642--1726/27)]{enlightenmentTeal}
        \small
        \textbf{Origin:} Woolsthorpe, England\\
        \textbf{Key work:} ``Method of fluxions'' (c.\,1665--1666) ---
        rates of change for physics.
      \end{personbox}

      \vspace{0.2em}
      \begin{personbox}[Gottfried Wilhelm Leibniz (1646--1716)]{enlightenmentTeal}
        \small
        \textbf{Origin:} Leipzig, Germany\\
        \textbf{Key work:} Independently developed calculus (published 1684--1686)
        with the $dy/dx$ notation we still use today.
      \end{personbox}
    \end{column}

    \begin{column}{0.49\textwidth}
      % -- Slide 19: Three derivative notations --
      \begin{tikzpicture}
        \node[draw=enlightenmentTeal!50, fill=enlightenmentTeal!5,
              rounded corners=2pt, inner sep=5pt, text width=1.4cm,
              align=center, font=\scriptsize]
          (fer) at (0,0) {\textbf{Fermat}\\[2pt]$\frac{f(x+e)-f(x)}{e}$\\[1pt]{\tiny $e\to 0$}};
        \node[draw=enlightenmentTeal!50, fill=enlightenmentTeal!5,
              rounded corners=2pt, inner sep=5pt, text width=1.4cm,
              align=center, font=\scriptsize]
          (new) at (2.3,0) {\textbf{Newton}\\[2pt]$\dot{y}$\\[1pt]{\tiny fluxion}};
        \node[draw=enlightenmentTeal!50, fill=enlightenmentTeal!5,
              rounded corners=2pt, inner sep=5pt, text width=1.4cm,
              align=center, font=\scriptsize]
          (lei) at (4.6,0) {\textbf{Leibniz}\\[2pt]$\dfrac{dy}{dx}$\\[1pt]{\tiny still used}};
        % Equals signs
        \node[font=\small] at (1.15,0) {$=$};
        \node[font=\small] at (3.45,0) {$=$};
        % Caption
        \node[font=\scriptsize\itshape, text=gray] at (2.3,-1.2)
          {Three notations, one idea: rate of change.};
      \end{tikzpicture}

      \vspace{0.3em}
      % -- Slide 19: Curve with tangent line --
      \begin{tikzpicture}[scale=0.8]
        \begin{axis}[
          width=5.5cm, height=4cm,
          axis lines=middle,
          ticks=none,
          xmin=-0.5, xmax=3.2, ymin=-0.3, ymax=4.5,
          clip=false
        ]
          % Curve f(x) = x^2/2
          \addplot[domain=0:3, smooth, thick, calculusRed, samples=40] {x*x/2};
          % Tangent line at x=1.5: f(1.5)=1.125, f'(1.5)=1.5
          \addplot[domain=0.3:2.7, thick, linalgTeal, dashed] {1.5*(x-1.5)+1.125};
          % Point of tangency
          \addplot[only marks, mark=*, mark size=2pt, calculusRed] coordinates {(1.5,1.125)};
          % Label
          \node[font=\tiny, text=linalgTeal] at (axis cs:2.5,2.8) {slope $= \frac{dy}{dx}$};
        \end{axis}
      \end{tikzpicture}

      \vspace{0.1em}
      {\scriptsize\itshape The derivative = the slope of this line.}

      \vspace{0.3em}
      {\small\textbf{LLM connection:} Training means computing
        derivatives of the error w.r.t.\ \emph{every} weight ---
        billions of weights, billions of derivatives, millions of times.}
    \end{column}
  \end{columns}

  \note{
    ``Newton and Leibniz famously fought over who invented calculus. Fermat was
    doing something similar decades earlier. None of them imagined that 300 years
    later, derivatives would be computed trillions of times to train a chatbot.''
    Note: Fermat also appears in Section 3 (Probability) with Pascal.
    When he comes up again, point out: ``Fermat again! He contributed to both
    calculus and probability --- two separate pillars of the LLM.''
  }
\end{frame}

% ----------------------------------------------------------
% Slide 20: The Chain Rule -- The Key to Backpropagation
% ----------------------------------------------------------
\begin{frame}{The Chain Rule --- The Key to Backpropagation}

  \begin{columns}[T]
    \begin{column}{0.47\textwidth}
      \textbf{Recipe} (the chain rule in plain language):

      \vspace{0.3em}
      If $y$ depends on $u$ and $u$ depends on $x$\ldots

      \vspace{0.2em}
      To find how much $x$ affected $y$ through the chain,
      \alert{multiply the effects at each step}.

      \vspace{0.5em}
      {\footnotesize\colorbox{enlightenmentTeal!8}{\parbox{0.92\linewidth}{%
        \textbf{Formal:}\quad
        $\dfrac{dy}{dx} = \dfrac{dy}{du} \cdot \dfrac{du}{dx}$\\[4pt]
        Leibniz's notation makes this look almost obvious ---
        the $du$'s ``cancel.'' (They don't really, but the intuition helps.)}}}

      \vspace{0.5em}
      \begin{insightbox}
        Backpropagation uses the chain rule to figure out which weights
        to blame for the error.
      \end{insightbox}
    \end{column}

    \begin{column}{0.50\textwidth}
      % -- Slide 20: Chain rule -- math + neural net --
      \begin{tikzpicture}[
        bx/.style={draw, rounded corners=2pt, inner sep=3pt,
                   font=\scriptsize, minimum width=0.8cm, minimum height=0.5cm},
        arr/.style={->,>=stealth,thick}
      ]
        % TOP ROW: Math version
        \node[font=\tiny\bfseries] at (0,3.2) {Math:};
        \node[bx, fill=gray!15] (x1) at (0,2.5) {$x$};
        \node[bx, fill=enlightenmentTeal!15] (f) at (1.8,2.5) {$f$};
        \node[bx, fill=gray!15] (u) at (3.2,2.5) {$u$};
        \node[bx, fill=enlightenmentTeal!15] (g) at (4.6,2.5) {$g$};
        \node[bx, fill=gray!15] (y1) at (6.0,2.5) {$y$};
        \draw[arr,probBlue] (x1)--(f);
        \draw[arr,probBlue] (f)--(u);
        \draw[arr,probBlue] (u)--(g);
        \draw[arr,probBlue] (g)--(y1);
        % BOTTOM ROW: Neural net version
        \node[font=\tiny\bfseries] at (0,1.0) {NN:};
        \node[bx, fill=gray!15] (xi) at (0,0.3) {Input};
        \node[bx, fill=linalgTeal!15] (l1) at (1.8,0.3) {Layer 1};
        \node[bx, fill=gray!15] (h) at (3.2,0.3) {Hidden};
        \node[bx, fill=linalgTeal!15] (l2) at (4.6,0.3) {Layer 2};
        \node[bx, fill=gray!15] (o) at (6.0,0.3) {Output};
        % Forward arrows (blue)
        \draw[arr,probBlue] (xi)--(l1);
        \draw[arr,probBlue] (l1)--(h);
        \draw[arr,probBlue] (h)--(l2);
        \draw[arr,probBlue] (l2)--(o);
        % Backward arrows (red)
        \draw[arr,calculusRed,dashed] (o.south) to[bend right=25] node[below,font=\tiny,text=calculusRed]{chain rule} (l2.south);
        \draw[arr,calculusRed,dashed] (l2.south) to[bend right=25] (h.south);
        \draw[arr,calculusRed,dashed] (h.south) to[bend right=25] (l1.south);
        \draw[arr,calculusRed,dashed] (l1.south) to[bend right=25] (xi.south);
      \end{tikzpicture}
    \end{column}
  \end{columns}

  \note{
    ``Backpropagation --- the algorithm that trains every neural network --- is
    `just' the chain rule applied over and over, layer by layer, backwards through
    the network. It's calculus class, running at scale.''
    Historical: the chain rule was implicit in Leibniz; made explicit by various
    18th--19th century mathematicians. No single ``inventor.''
    \verifytag{The precise attribution of the chain rule is diffuse --- Leibniz,
    Euler, and others all used it.}
  }
\end{frame}

% ----------------------------------------------------------
% Slide 21: Gradient Descent [EXTENDED]
% ----------------------------------------------------------
\begin{frame}{Gradient Descent --- Rolling Downhill \hfill {\small\textcolor{gray}{[EXTENDED]}}}

  \begin{columns}[T]
    \begin{column}{0.45\textwidth}
      \begin{personbox}[Augustin-Louis Cauchy (1789--1857)]{enlightenmentTeal}
        \small
        \textbf{Origin:} Paris, France\\
        \textbf{Key work:} Published the first description of the method of
        steepest descent (1847), \emph{M\'ethode g\'en\'erale pour la
        r\'esolution des syst\`emes d'\'equations simultan\'ees}.
      \end{personbox}

      \vspace{0.5em}
      \textbf{Recipe:}
      \begin{enumerate}
        \item Compute the gradient (``which way is uphill?'')
        \item Take a small step in the \alert{opposite} direction
        \item Repeat
      \end{enumerate}

      \vspace{0.3em}
      {\small Cauchy was solving equations, not training neural networks.\\
       But the algorithm is the same.}
    \end{column}

    \begin{column}{0.52\textwidth}
      % -- Slide 21: Gradient descent as 2D contour map --
      \begin{tikzpicture}[scale=0.8]
        % Contour ellipses (loss landscape top-down view)
        \draw[gray!30, thick] (2.5,1.5) ellipse(2.8 and 2.0);
        \draw[gray!40, thick] (2.5,1.5) ellipse(2.2 and 1.5);
        \draw[gray!50, thick] (2.5,1.5) ellipse(1.5 and 1.0);
        \draw[gray!60, thick] (2.5,1.5) ellipse(0.8 and 0.5);
        % Minimum label
        \node[circle, fill=calculusRed, inner sep=1.5pt] (min) at (2.5,1.5) {};
        \node[font=\tiny, below right, text=calculusRed] at (2.5,1.5) {minimum};
        % Gradient descent path (zigzag)
        \draw[->,thick,calculusRed]
          (0.2,3.3) -- (1.0,2.8) -- (1.8,2.5) -- (2.0,1.9) -- (2.3,1.7) -- (2.45,1.55);
        \node[circle, fill=gray, inner sep=1.5pt] at (0.2,3.3) {};
        \node[font=\tiny, above, text=gray] at (0.2,3.3) {start};
        % Labels
        \node[font=\tiny, text=gray] at (4.8,3.2) {high loss};
        \node[font=\tiny, text=calculusRed] at (3.5,1.0) {low loss};
        % Caption
        \node[font=\scriptsize\itshape, text=gray, text width=5cm, align=center]
          at (2.5,-0.8) {Gradient descent: always walk downhill.\\Cauchy, 1847.};
      \end{tikzpicture}
    \end{column}
  \end{columns}

  \note{
    This is an [EXTENDED] slide --- can be cut, covered verbally in the backprop slide.
    ``Cauchy described this in 1847 to solve equations. 170 years later, it's how
    every LLM on Earth is trained. The error is the height. The gradient tells you
    which way is downhill. You take a step. Repeat a billion times.''
    \verifytag{Cauchy 1847 paper on steepest descent --- confirm exact title and year.}
  }
\end{frame}

% ----------------------------------------------------------
% Slide 22: Backpropagation
% ----------------------------------------------------------
\begin{frame}{Backpropagation --- The Algorithm That Changed Everything}

  \begin{columns}[T]
    \begin{column}{0.50\textwidth}
      \begin{personbox}[{Seppo Linnainmaa (b.\,1945)}]{enlightenmentTeal}
        \small
        \textbf{Origin:} Finland\\
        \textbf{Key work:} Published automatic differentiation in his
        Master's thesis (1970, University of Helsinki) --- the mathematical
        core of backpropagation.
      \end{personbox}

      \vspace{0.2em}
      \begin{personbox}[Rumelhart/Hinton/Williams (1986)]{enlightenmentTeal}
        \small
        \textbf{David Rumelhart} (1942--2011, USA)\\
        \textbf{Geoffrey Hinton} (b.\,1947, London, UK/Canada)\\
        \textbf{Ronald Williams} (USA)\\[3pt]
        Published ``Learning representations by back-propagating errors''
        (1986, \emph{Nature}). Demonstrated backprop could train
        multi-layer networks. Reignited the field.\\[3pt]
        \textbf{Hinton shared the 2024 Nobel Prize in Physics}
        with John Hopfield for foundational work enabling machine learning.
      \end{personbox}
    \end{column}

    \begin{column}{0.47\textwidth}
      % -- Slide 22: Backpropagation network --
      \begin{tikzpicture}[
        neuron/.style={circle, draw, minimum size=10pt, inner sep=0pt, font=\tiny},
        scale=0.8
      ]
        % Input layer (3 nodes)
        \node[neuron, fill=gray!20] (i1) at (0,2.0) {};
        \node[neuron, fill=gray!20] (i2) at (0,1.0) {};
        \node[neuron, fill=gray!20] (i3) at (0,0.0) {};
        \node[font=\tiny,left] at (-0.3,1.0) {Input};
        % Hidden layer (4 nodes)
        \node[neuron, fill=linalgTeal!20] (h1) at (2,2.5) {};
        \node[neuron, fill=linalgTeal!20] (h2) at (2,1.5) {};
        \node[neuron, fill=linalgTeal!20] (h3) at (2,0.5) {};
        \node[neuron, fill=linalgTeal!20] (h4) at (2,-0.5) {};
        \node[font=\tiny] at (2,-1.1) {Hidden};
        % Output layer (2 nodes)
        \node[neuron, fill=accentOrange!20] (o1) at (4,1.5) {};
        \node[neuron, fill=accentOrange!20] (o2) at (4,0.5) {};
        \node[font=\tiny,right] at (4.3,1.0) {Output};
        % Forward connections (blue, thin)
        \draw[->,thin,probBlue!60] (i1)--(h1); \draw[->,thin,probBlue!60] (i1)--(h2);
        \draw[->,thin,probBlue!60] (i1)--(h3); \draw[->,thin,probBlue!60] (i1)--(h4);
        \draw[->,thin,probBlue!60] (i2)--(h1); \draw[->,thin,probBlue!60] (i2)--(h2);
        \draw[->,thin,probBlue!60] (i2)--(h3); \draw[->,thin,probBlue!60] (i2)--(h4);
        \draw[->,thin,probBlue!60] (i3)--(h1); \draw[->,thin,probBlue!60] (i3)--(h2);
        \draw[->,thin,probBlue!60] (i3)--(h3); \draw[->,thin,probBlue!60] (i3)--(h4);
        \draw[->,thin,probBlue!60] (h1)--(o1); \draw[->,thin,probBlue!60] (h1)--(o2);
        \draw[->,thin,probBlue!60] (h2)--(o1); \draw[->,thin,probBlue!60] (h2)--(o2);
        \draw[->,thin,probBlue!60] (h3)--(o1); \draw[->,thin,probBlue!60] (h3)--(o2);
        \draw[->,thin,probBlue!60] (h4)--(o1); \draw[->,thin,probBlue!60] (h4)--(o2);
        % Forward label
        \draw[->,thick,probBlue] (0.5,3.0) -- (3.5,3.0)
          node[midway,above,font=\tiny\bfseries] {Forward pass};
        % Backward label
        \draw[->,thick,calculusRed,dashed] (3.5,-1.5) -- (0.5,-1.5)
          node[midway,below,font=\tiny\bfseries] {Backward pass (chain rule)};
      \end{tikzpicture}

      \vspace{0.3em}
      {\small\textbf{Timeline:}}

      \vspace{0.2em}
      {\footnotesize
      \begin{tabular}{@{}rl@{}}
        1847 & Cauchy --- steepest descent \\
        1970 & Linnainmaa --- auto.\ differentiation \\
        1986 & Rumelhart/Hinton/Williams --- backprop \\
        2024 & Hinton --- Nobel Prize in Physics \\
      \end{tabular}}
    \end{column}
  \end{columns}

  \note{
    ``Linnainmaa figured out the algorithm in 1970. Rumelhart, Hinton, and Williams
    showed it could train neural networks in 1986. Hinton shared the Nobel Prize
    in Physics in 2024 with John Hopfield for foundational work enabling machine
    learning with artificial neural networks.''
    \verifytag{Linnainmaa's 1970 thesis --- some sources describe it as reverse-mode
    automatic differentiation. Others credit earlier related ideas (Bryson, Dreyfus,
    Kelley in the 1960s, control theory). Present Linnainmaa as key formalization
    but note broader lineage.}
    \verifytag{Nobel Prize 2024 --- Hinton and Hopfield shared the Nobel Prize
    in Physics 2024. Confirm exact citation wording.}
  }
\end{frame}

% ----------------------------------------------------------
% Slide 23: Stochastic Gradient Descent
% ----------------------------------------------------------
\begin{frame}{Stochastic Gradient Descent --- Learning from Samples}

  \begin{columns}[T]
    \begin{column}{0.45\textwidth}
      \begin{personbox}[Robbins \& Monro (1951)]{enlightenmentTeal}
        \small
        \textbf{Herbert Robbins} (1915--2001, New Castle, PA, USA)\\
        \textbf{Sutton Monro} (1920--1995, USA)\\[3pt]
        Published the Robbins--Monro algorithm (1951) --- the first
        stochastic approximation method.
        The theoretical foundation for SGD.
      \end{personbox}

      \vspace{0.5em}
      \begin{insightbox}
        Every LLM you have ever used was trained with some variant of SGD.
        The Robbins--Monro paper from 1951 is running inside every AI lab today.
      \end{insightbox}
    \end{column}

    \begin{column}{0.52\textwidth}
      % -- Slide 23: Full GD vs SGD comparison --
      \begin{tikzpicture}[scale=0.7]
        % LEFT: Full GD
        \begin{scope}[xshift=0cm]
          \node[font=\scriptsize\bfseries] at (2,3.5) {Full GD};
          % Contours
          \draw[gray!40] (2,1.5) ellipse(2.0 and 1.5);
          \draw[gray!50] (2,1.5) ellipse(1.3 and 0.9);
          \draw[gray!60] (2,1.5) ellipse(0.5 and 0.3);
          % Smooth path
          \draw[->,thick,probBlue] (0.3,3.0) -- (1.0,2.5) -- (1.5,2.0) -- (1.9,1.6);
          \node[circle,fill=gray,inner sep=1pt] at (0.3,3.0) {};
          \node[circle,fill=calculusRed,inner sep=1.5pt] at (2,1.5) {};
          \node[font=\tiny\itshape, text=gray, text width=2.5cm, align=center]
            at (2,-0.5) {ALL data per step.\\Smooth but slow.};
        \end{scope}
        % RIGHT: SGD
        \begin{scope}[xshift=5.5cm]
          \node[font=\scriptsize\bfseries] at (2,3.5) {SGD};
          % Contours
          \draw[gray!40] (2,1.5) ellipse(2.0 and 1.5);
          \draw[gray!50] (2,1.5) ellipse(1.3 and 0.9);
          \draw[gray!60] (2,1.5) ellipse(0.5 and 0.3);
          % Noisy zigzag path
          \draw[->,thick,calculusRed]
            (0.3,3.0) -- (1.2,2.3) -- (0.8,2.6) -- (1.5,2.0)
            -- (1.2,1.8) -- (1.8,1.6) -- (1.6,1.7) -- (1.95,1.5);
          \node[circle,fill=gray,inner sep=1pt] at (0.3,3.0) {};
          \node[circle,fill=calculusRed,inner sep=1.5pt] at (2,1.5) {};
          \node[font=\tiny\itshape, text=gray, text width=2.5cm, align=center]
            at (2,-0.5) {Random sample per step.\\Noisy but fast.};
        \end{scope}
      \end{tikzpicture}
    \end{column}
  \end{columns}

  \note{
    ``Training GPT-4 used trillions of tokens. You can't compute the gradient
    on all of them at once. SGD says: just grab a random batch, compute the
    gradient on that, take a step, repeat. It's noisy, but it works. And it's fast.''
    LLM connection: every LLM ever deployed was trained with some variant of SGD.
    \verifytag{Monro's dates (1920--1995) and birthplace --- verify.}
  }
\end{frame}

% ----------------------------------------------------------
% Slide 24: The Loss Landscape -- A Visual Tour
% ----------------------------------------------------------
\begin{frame}{The Loss Landscape --- A Visual Tour}

  \begin{center}
    \visualplaceholder[12cm]{3D LOSS LANDSCAPE (to be generated as figures/loss-landscape.pdf via Python/matplotlib or pgfplots standalone).
      Style: Li et al., 2018 ``Visualizing the Loss Landscape of Neural Nets.''
      Surface: 3D plot of $f(x,y) = 0.5(x^2+y^2) + 3\exp(-2((x-1)^2+(y+1)^2)) - \exp(-3((x+1.5)^2+(y-0.5)^2))$
      or similar multi-modal surface.
      Labeled annotations:
      (A) ``Global minimum'' -- deepest valley, green marker
      (B) ``Local minimum'' -- secondary valley, yellow marker
      (C) ``Saddle point'' -- flat ridge, gray marker
      (D) ``SGD path'' -- noisy red trajectory from random start to global min
      Colormap: viridis or coolwarm. View angle: 30\textdegree{} elevation.}
  \end{center}

  \vspace{-0.3em}
  {\small\centering\textcolor{gray}{This is what training an LLM looks like
    mathematically. The landscape has \emph{billions} of dimensions ---
    we can only visualize a 2D slice.}\par}

  \note{
    ``This is what training an LLM looks like mathematically. You start at a random
    point on this landscape and try to find the lowest valley. The landscape has
    billions of dimensions --- we can only visualize a 2D slice.''
    The ``loss'' measures how wrong the network's predictions are.
    Training = finding the lowest point on this landscape.
    SGD's noisy path actually helps --- it can escape shallow local minima
    that full gradient descent might get stuck in.
  }
\end{frame}

% ----------------------------------------------------------
% Slide 25: Convergence Diagram Update -- Calculus & Optimization
% ----------------------------------------------------------
\begin{frame}{Building the Diagram --- Calculus Complete}

  \begin{columns}[T]
    \begin{column}{0.55\textwidth}
      \begin{center}
        \visualplaceholder[6.5cm]{CONVERGENCE DIAGRAM -- VARIANT 2 (figures/convergence-2.pdf).
          Node (1) ``Linear Algebra'' lit in linalgTeal (from Section 1).
          Node (2) ``Calculus \& Optimization'' now lit in calculusRed.
          Label: ``Gradient descent = how it learns.''
          Nodes (3)--(10) remain gray with dashed arrows.
          Central node still ``?'' in gray.}
      \end{center}
    \end{column}

    \begin{column}{0.42\textwidth}
      \textbf{Summary --- who taught it to learn:}

      \vspace{0.3em}
      {\small
      \begin{tabular}{@{}rl@{}}
        c.\,250 BCE & Archimedes --- proto-calculus \\
        c.\,1636 & Fermat --- tangents, extrema \\
        1665--66 & Newton --- fluxions \\
        1684--86 & Leibniz --- $dy/dx$ notation \\
        1847 & Cauchy --- steepest descent \\
        1951 & Robbins--Monro --- SGD \\
        1970 & Linnainmaa --- auto.\ diff. \\
        1986 & Rumelhart/Hinton/Williams \\
               & \quad --- backpropagation \\
      \end{tabular}}

      \vspace{0.5em}
      \begin{insightbox}
        Gradient descent + backpropagation = the learning algorithm of every LLM.
      \end{insightbox}
    \end{column}
  \end{columns}

  \note{
    Transition: ``The network can now learn. But learn WHAT? It learns to predict
    probability distributions over words. For that, we need probability theory.''
    Connection arrow: ``Gradient descent + backpropagation = the learning algorithm
    of every LLM.''
    This is the end of Session 1 (Slides 1--25). Natural break point.
    ``We now know the skeleton of a neural network and how it learns.
    Next: how does it deal with uncertainty?''
  }
\end{frame}